\paragraph{Multi-agent communication.}
Agent communication languages~\citep{finin_kqml_1994} defined formal message types, and learned communication~\citep{sukhbaatar_commnet_2016, foerster_dial_2016} later optimized bandwidth end-to-end. Recent work examines when communication helps~\citep{rizvi_martel_communication_2026}, what to communicate~\citep{chen_five_ws_2026}, and why agents communicate~\citep{zhou_why_agents_communicate_2025}; a communication-centric survey~\citep{yan_beyond_self_talk_2025} organizes this landscape. Multi-agent debate~\citep{du_debate_2024}, orchestration frameworks~\citep{li_camel_2023, hong_metagpt_2024, wu_autogen_2024, chen_agentverse_2024, wang_megaagent_2025}, and information propagation analysis~\citep{shen_eib_2025} confirm that communication structure affects performance. Topology learning~\citep{li_arg_designer_2026, lu_dytopo_2026} automates structure design. These contributions operate at the system level; we work at the message level, testing whether functional categories predict which messages are safe to prune.

\paragraph{Communication efficiency.}
Existing methods operate at granularities coarser than message function. Agent-level pruning~\citep{wang_agentdropout_2025, wang_agentdropout_v2_2026, li_agp_2025, xiao_agentdiet_2026} removes entire agents. Instance-level routing~\citep{wang_dagp_2025, fan_imad_2026} predicts which queries require multi-agent collaboration. Turn-level stopping~\citep{zeng_s2mad_2025} halts conversations early. Topology optimization~\citep{chen_optima_2025, chen_goagent_2026, wang_anymac_2025, li_lacp_2025} restructures agent organization. Token-level filtering~\citep{safesieve_2025, lin_supervisoragent_2026} selectively removes content via safety/relevance filters. KVCOMM~\citep{ye_kvcomm_2025} shares cached key--value states across contexts (7.8$\times$ speedup). Latent-space communication~\citep{zou_latentmas_2025, du_interlat_2025} and holistic optimization~\citep{li_omac_2025, jiang_agentcomm_2026} pursue complementary directions (Appendix Table~\ref{tab:method_comparison}). Debate-level analyses suggest most accuracy gains are attributable to voting rather than information exchange~\citep{choi_debate_or_vote_2025, wu_can_agents_debate_2025}; distillation internalizes debate into single models~\citep{yi_latent_agents_2026}. None of these methods link communicative function to pruning outcomes; our taxonomy-driven pruning identifies safe categories based on functional role.

\paragraph{Structural versus functional pruning.}
AgentPrune~\citep{zhang_agentprune_2025} selects edges to retain in the communication graph, deciding \emph{who talks to whom}. Our method addresses an orthogonal question: given a fixed topology, \emph{what} does each message contribute? The two operate at different granularities (edge vs.\ message; Table~\ref{tab:method_comparison}) and compose in principle, though empirical validation remains future work.

\paragraph{Theoretical foundations.}
Our taxonomy draws on speech act theory~\citep{austin_how_1962, searle_speech_acts_1969} and dialogue act schemes (DAMSL~\citep{core_damsl_1997}, ISO 24617-2~\citep{bunt_iso_2012}, Switchboard DA~\citep{stolcke_dialogue_acts_2000}). Backchannels~\citep{yngve_getting_1970}, continuers~\citep{schegloff_discourse_1982}, and acknowledgment tokens~\citep{jefferson_acknowledgment_1984} carry minimal propositional content yet serve a load-bearing role in conversational grounding~\citep{clark_contributing_1989}. Computational studies~\citep{inoue_backchannel_2025, qian_backchannel_2026} confirm these signals persist in neural systems; LLM extensions~\citep{qamar_dialogue_acts_2025, zhang_communication_strategies_2025} show standard categories apply to agent conversations. Attribution methods~\citep{xia_hivemind_2026, horovicz_agentshap_2025, cemri_mast_2025} assign credit at the agent level but do not decompose messages by function. We extend these foundations by connecting functional categories to pruning outcomes through causal intervention.
